% ======================================================================
% SECTION: Introduction
% ======================================================================
\section{Introduction}
\label{sec:intro}
Expressive 3D human avatars, represented using parametric body models such as SMPL-X~\cite{smplx}, have broad applications in interactive entertainment, embodied conversational agents, and immersive telepresence.
However, automatically generating full-body avatar animation conditioned on speech remains a fundamentally difficult problem.
Beyond generation itself, producing motion that appears natural is equally challenging: the face, hands, and body must move in precise synchrony with the rhythm, prosody, and meaning of speech.
Achieving this coherent, speech-driven full-body coordination is a long-standing open problem in computer graphics and animation~\cite{mcneill1992hand,s2g,beat,emage2024}.
Progress in generating 3D human avatars conditioned on speech has grown with the release of large scale motion capture datasets such as BEAT and BEAT2~\cite{beat,emage2024}.
However, despite this progress, several limitations remain.

First, recent methods~\cite{talkshow2023,gesturediffuclip,syntalker} rely on Vector Quantized Variational Autoencoders (VQ-VAE)~\cite{vqvae,vqvae2} to discretize continuous motion into codebook tokens.
While this simplifies modeling, it introduces quantization artifacts.
Since human motion is continuous, mapping it to discrete tokens can cause abrupt transitions when switching between codewords, resulting in visible joint snapping and reduced smoothness.
Second, generating long sequences of avatar motion remains a challenge.
Prior models are trained on short, fixed length clips and do not naturally generalize to extended recordings.
Existing strategies such as seed frame conditioning or post hoc smoothing attempt to mitigate boundary artifacts, but they do not explicitly train the model to enforce kinematic continuity across segment boundaries.
Third, previous methods depend on gesture type labels provided by specific datasets such as BEAT2.
This dependence limits applicability to datasets without such annotations.
Since manual labeling is expensive and time consuming, an automated method for generating pseudo labels is desirable.
Lastly, although BEAT2 contains recordings from 25 speakers, previous studies predominantly report performance on a single speaker, typically Speaker~2~\cite{emage2024,syntalker}.
This raises concerns about the generalizability of existing models.

To address these gaps, we present \textbf{GestureCraft}, a speech to motion framework.
Our contributions are summarized as follows:

\begin{itemize}[leftmargin=*]
    \item \textbf{Hierarchical VAE.} In our framework, the VAE operates in a continuous latent space, which avoids quantization artifacts. We introduce a hierarchical design in which the latent representation is partitioned into face, body, and hand components. Each component is trained with a partial reconstruction objective, encouraging specialization and improving local motion quality. In addition, a global reconstruction objective ensures synchronization and coherence across all parts of the body.
    \item \textbf{Anchor-conditioned decoder.} We introduce anchor prefix conditioning to enable smooth chunked generation, where each segment is conditioned on the final frames of the previous chunk to preserve temporal continuity. To further improve boundary consistency, we propose a continuation velocity loss that explicitly penalizes velocity discrepancies across chunk boundaries.
    \item \textbf{LLM-based gesture type planner.}
    We introduce a finetuned Large Language Model (LLM) that predicts BEAT2 style gesture type annotations from raw speech transcripts, i.e., the planner generates temporal segments that specify when beat or semantic gestures should occur. By generating these annotations automatically, our framework preserves downstream motion quality while removing the need for gesture labels during inference.
    \item \textbf{Cross speaker generalization.} Existing studies evaluate only on Speaker~2 in the BEAT2 dataset. In contrast, we evaluate our model on the full 25-speakers BEAT2 test set to assess its robustness and generalization across speakers.
\end{itemize}

% ======================================================================
% SECTION: Related Work
% ======================================================================

\section{Related Work}
\label{sec:related}
\noindent \textbf{Co-speech gesture generation.}
% \citet{mcneill1992hand} established a foundational taxonomy of co speech gestures by distinguishing between beat gestures and semantic gestures. Beat gestures are rhythmic movements that align with speech prosody and provide emphasis without conveying semantic meaning. In contrast, semantic gestures, including iconic and metaphoric forms, visually represent the content of speech.
Early approaches to 3D human avatar motion generation learn mappings from speech features to skeletal joint rotations using RNNs or CNNs~\cite{s2g,trimodal,ha2g,habibie}, with most methods restricted to motion only of the upper body or the hand.
To address this limitation, more recent work has extended these methods to full-body SMPL-X generation and adopted more modern architectures.
TalkSHOW~\cite{talkshow2023} generates holistic body motion from speech using a compositional VQ-VAE for body and hand motion together with a Gated PixelCNN autoregressive prior over codebook indices.
While this model can produce long form motion, its quality depends on rollout stability since it does not explicitly enforce boundary continuity.
GestureDiffuCLIP~\cite{gesturediffuclip} runs latent diffusion in a VQ-VAE pretrained gesture latent space with CLIP-based style control, demonstrating that VQ-then-diffusion pipelines can produce stylized co-speech motion.
DiffuseStyleGesture~\cite{diffusestylegesture2023} instead applies diffusion directly in the continuous pose space with classifier-free guidance.
However, long sequence generation still relies on initial frame conditioning rather than a learned mechanism for enforcing continuity across temporal windows.
Furthermore, all prior methods evaluated on BEAT2 report results only on a single speaker (Speaker~2), leaving generalization across the full 25-speakers set largely unverified.

\noindent \textbf{Hierarchical latent variable models.}
VQ-VAE~\cite{vqvae} is a latent variable model that encodes continuous inputs into a sequence of discrete codebook indices, replacing the continuous latent space with a finite set of learned embeddings.
The decoder reconstructs the data from these discrete tokens, while generation is typically performed using an autoregressive prior over the code indices.
VQ-VAE 2~\cite{vqvae2} extends this formulation by introducing a multi-scale hierarchy of codebooks operating at different resolutions, where separate latent grids represent coarse and fine spatial structure.
This results in a layered discrete representation rather than a single flat codebook.
\citet{adiban2022} further extend this direction with hierarchical residual vector quantization (HR-VQVAE), where multiple codebook levels are applied sequentially and each level quantizes the residual error of the previous one.
This produces a hierarchical discrete representation with multiple stages of refinement, at the cost of increased decoding complexity due to sequential codebook lookup across levels.
% HR-VQVAE introduces discrete codebooks that are effective for modeling image pixels.
However, in 3D human avatar motion generation, SMPL-X pose parameters lie in a continuous space, making direct application of HR-VQVAE less suitable.
We therefore retain the hierarchical coarse-to-fine inductive bias of HR-VQVAE, but replace discrete quantization with a continuous Gaussian latent.
This latent is partitioned channel wise, with each slice supervised on a specific body part, namely, the face, body, and hands.

\noindent \textbf{Flow matching and diffusion priors over latents.}
Denoising diffusion models~\cite{ddpm} and their extensions, including latent diffusion and flow matching, provide powerful frameworks for modeling continuous data distributions.
Flow matching~\cite{lipman2023flow} formulates generative modeling as learning a continuous velocity field without simulating the full diffusion process, while rectified flow~\cite{rectifiedflow} further simplifies sampling by encouraging near-linear transport paths that can be integrated in a few ODE steps.
We adopt the rectified flow formulation to learn a Transformer~\cite{transformer} based velocity field over our temporal latent sequence, conditioned on speaker speech, identity, and gesture type.
% Instead of applying diffusion directly in SMPL X pose space, we operate in a compressed latent space (temporal stride 8, latent dimension 16), which reduces the computational cost of each ODE step and separates the complexity of motion generation from output resolution.

\noindent \textbf{LLM-driven motion planning.}
Recent work has explored the use of LLMs for motion generation tasks.
For example, MECo~\cite{meco2025} finetunes an LLM to process audio together with motion examples and directly generates gesture tokens as output.
In contrast, our approach uses an LLM for gesture type planning.
It produces timestamped segments with predicted gesture types, such as beat and semantic.
This structured temporal plan is then used to condition the downstream VAE for motion generation.

% ======================================================================
% SECTION: Method
% ======================================================================

\section{Method}
\label{sec:method}
\begin{figure*}[t]
  \centering
  \includegraphics[width=\textwidth]{system_overview.png}

  \Description{Three-band system overview. (a) Planning: a transcript is
    fed into a LoRA-fine-tuned Qwen language model that emits a sequence
    of timed gesture-type segments. (b) Generation: audio, speaker, and
    gesture-type embeddings condition a flow-matching prior that samples
    a temporal latent partitioned channel-wise into face, body, and hand
    slices; the latent and an anchor-prefix from the previous chunk drive
    a Transformer decoder with dual cross-attention (to the body-part-hierarchical latent and to audio), whose
    output is rendered through the SMPL-X body model. (c) Long-form: the
    recording is split into successive chunks; the last K frames of each
    chunk are reused as the anchor prefix of the next, and a
    continuation velocity loss penalizes velocity discontinuity at chunk
    boundaries.}
  \caption{\textbf{GestureCraft inference pipeline.}
    (1) Speech and transcript $\rightarrow$
    (2) LLM planner emits beat/semantic gesture-type segments
    (Section~\ref{sec:planner}) $\rightarrow$
    (3) per-chunk conditioning (audio, speaker, gesture type) $\rightarrow$
    (4) flow-matching prior samples temporal latent $\mathbf{z}$
    (Section~\ref{sec:prior}) $\rightarrow$
    (5) channel-wise face/body/hands partition
    $\mathbf{z}^{(0)},\mathbf{z}^{(1)},\mathbf{z}^{(2)}$
    (Section~\ref{sec:vae}) $\rightarrow$
    (6) hierarchical Transformer decoder $\rightarrow$
    (7) SMPL-X pose at $30$\,fps.
    Bottom: chunked autoregression reuses the previous chunk's last $K$
    frames as an anchor prefix
    (Section~\ref{sec:autoreg}).}
  \label{fig:overview}
\end{figure*}

We propose a method for co-speech gesture synthesis that produces continuous full-body SMPL-X motions conditioned on speech audio, speaker identity, and gesture type.
Our model, illustrated in Figure~\ref{fig:overview}, consists of four components: (i) Hierarchical continuous temporal VAE for structured motion representation, (ii) Conditional flow-matching prior that samples the latent conditioned on speaker speech, identity, and gesture type, (iii) Anchor-conditioned chunk-wise autoregressive decoder for long-form generation, and (iv) LLM-based gesture type planner for temporal segmentation.

\noindent \textbf{Notation.}
Let $\mathbf{x} \in \mathbb{R}^{T \times 165}$ denote a motion sequence at 30\,fps in an SMPL-X axis-angle representation.
The pose vector is decomposed into body, face, and hands,
$\mathbf{x}_\text{B} \in \mathbb{R}^{T \times 66}$, $\mathbf{x}_\text{F} \in \mathbb{R}^{T \times 9}$, and $\mathbf{x}_\text{H} \in \mathbb{R}^{T \times 90}$.
Speech is represented using Wav2Vec 2.0~\cite{wav2vec2} BASE features
$\mathbf{a} \in \mathbb{R}^{T \times 768}$.
Note that $\mathbf{x}$ encodes joint rotations only; root translation and global orientation drift are not modeled.
Each motion segment is associated with a gesture type label $g \in \{\text{beat}, \text{semantic}\}$ and a speaker identity $s$.

% -----------------------------------------------------------------------
\subsection{Hierarchical VAE}
\label{sec:vae}
% -----------------------------------------------------------------------
We model motion sequences using a variational autoencoder that encodes input motion into a continuous Gaussian latent sequence.
The latent is temporally downsampled and factorized into structured channel-wise groups corresponding to face, body, and hands.

\noindent \textbf{Encoding.}
Given motion sequence $\mathbf{x}$ and audio $\mathbf{a}$, we embed both using a Transformer encoder with positional encoding.
Audio features are added to the motion embeddings at the input level.
The encoder output is passed through a strided 1D convolution with stride $r=8$, producing a latent sequence $\mathbf{h} \in \mathbb{R}^{T' \times d}$.
Each latent token is projected to Gaussian parameters $(\boldsymbol{\mu}, \boldsymbol{\sigma})$, and the latent $\mathbf{z}$ is sampled via the reparameterization trick.

\noindent \textbf{Hierarchical latent factorization.}
The latent is split along the channel dimension into three groups $\mathbf{z} = [\mathbf{z}^{(0)}, \mathbf{z}^{(1)}, \mathbf{z}^{(2)}]$,
corresponding to face, body, and hands.
The latent dimension $d_z=16$ is partitioned as $(5,5,6)$ across the three groups.

\noindent \textbf{Decoding.}
A Transformer decoder reconstructs motion at full temporal resolution.
Each layer applies self-attention over motion tokens and cross-attention over two conditioning sources: (i) Latent tokens $\mathbf{z}$, augmented with identity and gesture type embeddings and (ii) Audio features.
This design allows the decoder to jointly leverage identity and timing information.

\noindent \textbf{Reconstruction objective.}
We employ a hierarchical reconstruction scheme with partial supervision over body parts.
The decoder produces intermediate outputs using progressively more latent slices.
The face is reconstructed from $\mathbf{z}^{(0)}$, the body from $[\mathbf{z}^{(0)}, \mathbf{z}^{(1)}]$, and the hands from all slices.
The corresponding loss is as follows:

\begin{equation}
\mathcal{L}^{(\ell)}_\text{local}
= \mathcal{L}_\text{aa}^{(\ell)}
+ \lambda_\text{vel}\mathcal{L}_\text{vel}^{(\ell)}
+ \mathcal{L}_\text{mesh}^{(\ell)}
+ \mathcal{L}_\text{joint}^{(\ell)}.
\end{equation}
The global reconstruction loss applied to the full output is:
\begin{equation}
\mathcal{L}_\text{global}
= \mathcal{L}_\text{aa}
+ \lambda_\text{vel}\mathcal{L}_\text{vel}
+ 5\mathcal{L}_\text{mesh}
+ 3\mathcal{L}_\text{joint}.
\end{equation}
The slice-wise KL regularization is:
\begin{equation}
\mathcal{L}_\text{KL}
= \sum_{\ell=0}^{2}
\mathrm{KL}(q(\mathbf{z}^{(\ell)}) \| \mathcal{N}(0, I)).
\label{eq:kl}
\end{equation}

\begin{figure*}[t]
  \centering
  \includegraphics[width=\textwidth]{vae_training.png}
  \Description{Diagram of the VAE training architecture. The encoder
    is a Transformer followed by a strided 1-D convolution that
    downsamples the input pose sequence by a factor of 8. The output
    is projected to mean and log-variance vectors, and the latent z
    is sampled via the reparameterization trick. Each latent token is
    partitioned channel-wise into three slices, colour-coded for face,
    body, and hands, which feed the decoder. The decoder reconstructs
    the SMPL-X pose sequence; per-level losses are applied separately
    to the face, body, and hand joints, in addition to a whole-body
    global loss.}
  \caption{VAE training architecture. The continuous latent $\mathbf{z}$ is partitioned channel-wise into three ordered slices (face/body/hands) and supervised by per-level body-part-localized reconstruction losses in addition to a whole-body global loss.}
  \label{fig:vae_train}
\end{figure*}

% -----------------------------------------------------------------------
\subsection{Conditional Flow-Matching Prior}
\label{sec:prior}
% -----------------------------------------------------------------------
We model the latent distribution of motion representations using a conditional flow-matching model.
Instead of sampling from a fixed Gaussian prior, we learn a velocity field $p(\mathbf{z} \mid \mathbf{a}, s, g)$ to generate latents conditioned on speech audio $\mathbf{a}$, speaker identity $s$, and gesture type $g$.

\noindent \textbf{Conditioning architecture.}
We construct a temporal conditioning representation using an audio encoder.
Wav2Vec 2.0~\cite{wav2vec2} features are first projected to a shared embedding space and combined with identity and gesture type embeddings.
When available, anchor-prefix embeddings are also injected into the same representation.
The resulting sequence is processed by a Transformer with sinusoidal positional encoding, producing a conditioning context $\mathbf{c}_{1:T'} \in \mathbb{R}^{T' \times d_\text{model}}$.
This context is used to condition the generative process of the latent sequence.

\noindent \textbf{Flow-matching model.}
We learn a neural velocity field $v_\theta(\mathbf{z}_t, t, \mathbf{c})$ that transforms Gaussian noise into structured latent codes.
The model is conditioned on both the intermediate noisy latent $\mathbf{z}_t$ and the conditioning context $\mathbf{c}$.
Time information $t \in [0,1]$ is injected using AdaLN modulation~\cite{dit}, while cross-attention is used to integrate audio and conditioning signals.

\noindent \emph{Stage 1: Flow-matching pretraining.}
We first train the model to match the VAE encoder posterior.
Given a latent sample $\mathbf{z}^* \sim q_\phi(\mathbf{z} \mid \mathbf{x}, \mathbf{a}, s, g)$, we construct a noisy interpolation:

\begin{equation}
    \mathbf{z}_t = (1 - t)\boldsymbol{\varepsilon} + t\mathbf{z}^*, \quad \boldsymbol{\varepsilon} \sim \mathcal{N}(\mathbf{0}, \mathbf{I}),
\end{equation}
and optimize the rectified-flow objective:

\begin{equation}
    \mathcal{L}_\text{flow} =
    \mathbb{E}_{t, \boldsymbol{\varepsilon}}
    \left\| v_\theta(\mathbf{z}_t, t, \mathbf{c}) - (\mathbf{z}^* - \boldsymbol{\varepsilon}) \right\|^2.
\end{equation}

\noindent \emph{Stage 2: Perceptual fine-tuning.}
We further refine the flow model using perceptual supervision through a frozen VAE decoder.
This allows optimization directly in motion space while keeping the latent generator unchanged.
The objective combines flow matching with motion-quality losses:

\begin{equation}
\mathcal{L}_\text{e2e} =
\mathcal{L}_\text{flow}
+ \lambda_\text{fgd} \mathcal{L}_\text{FGD}
+ \lambda_\text{beat} \mathcal{L}_\text{beat}.
\end{equation}
The FGD loss matches feature statistics between generated and ground-truth motion, while the beat loss encourages alignment between motion peaks and audio rhythm.
To enable differentiable optimization, we estimate a one-step latent prediction:

\begin{equation}
    \hat{\mathbf{z}}^* = \mathbf{z}_t + (1 - t)\, v_\theta(\mathbf{z}_t, t, \mathbf{c}),
\end{equation}
which is decoded through the frozen VAE during training.

\noindent \emph{Sampling.}
At inference time, we generate latents by integrating the learned velocity field from Gaussian noise using a midpoint ODE solver with $N=50$ steps.
The resulting latent sequence is decoded using the VAE decoder to produce motion.
A classifier-free guidance scale $w$ controls the trade-off between fidelity and diversity.

% -----------------------------------------------------------------------
\subsection{Anchor-Conditioned Decoder}
\label{sec:autoreg}
% -----------------------------------------------------------------------
Long motion sequences cannot be generated in a single forward pass due to memory constraints and sequence length limitations.
However, independently generated chunks often produce discontinuities at segment boundaries.
We address this using anchor-conditioned decoding during both training and inference, combined with chunk-wise autoregressive generation.

\noindent \textbf{Anchor-prefix conditioning.}
At each decoding step, the model receives an anchor-prefix consisting of the last $K$ frames from the previously generated chunk, denoted as $\mathbf{x}_a \in \mathbb{R}^{K \times 165}$, together with the corresponding audio segment $\mathbf{a}_a \in \mathbb{R}^{K \times 768}$.
The anchor motion is appended to the decoder input sequence, allowing self-attention to access prior motion context.
The anchor audio is similarly appended to the audio memory used in cross-attention, ensuring temporal alignment between motion and speech context.
During training, the decoder also reconstructs the anchor frames as an auxiliary objective to stabilize boundary prediction.

\noindent \textbf{Anchor curriculum.}
To reduce the discrepancy between training and inference, we apply a curriculum over anchor length $K$.
Specifically, we sample:

\begin{itemize}[leftmargin=*]
    \item 60\% full anchors ($K = K_\text{max}$), matching inference conditions.
    \item 20\% no anchors ($K = 0$), enabling fallback robustness.
    \item 20\% random anchors ($K \in [1, K_\text{max}-1]$), improving generalization to partial context.
\end{itemize}

\noindent \textbf{Continuation velocity loss.}
To improve temporal consistency at chunk boundaries, we introduce a weighted velocity matching loss to emphasize the transition region between anchor and generated motion.
Let $\hat{\mathbf{x}}_i$ and $\mathbf{x}_i$ denote predicted and ground-truth frames, and $\mathbf{x}_a^{(K)}$ the last anchor frame.
We define velocity differences:

\begin{align}
\hat{\Delta}_i &=
\begin{cases}
\hat{\mathbf{x}}_0 - \mathbf{x}_a^{(K)}, & i = 0 \\
\hat{\mathbf{x}}_i - \hat{\mathbf{x}}_{i-1}, & i > 0
\end{cases}, \\
\Delta_i &=
\begin{cases}
\mathbf{x}_0 - \mathbf{x}_a^{(K)}, & i = 0 \\
\mathbf{x}_i - \mathbf{x}_{i-1}, & i > 0.
\end{cases}
\end{align}
The continuation loss is computed over a horizon of $N$ frames:

\begin{equation}
\mathcal{L}_\text{cont} =
\frac{1}{\sum_{i=0}^{n-1} w_i}
\sum_{i=0}^{n-1} w_i \|\hat{\Delta}_i - \Delta_i\|^2,
\quad w_i = 1 - \frac{i}{N},
\label{eq:cont}
\end{equation}
where $n = \min(N, T)$.
This weighting emphasizes consistency at the anchor boundary while gradually relaxing constraints deeper into the generated chunk.

\noindent \textbf{Chunk-wise autoregressive generation.}
At inference time, motion is generated sequentially following the segmentation produced by the gesture type planner (Section~\ref{sec:planner}).
Each segment is further divided into fixed-length chunks of up to 150 frames.
Generation proceeds autoregressively as follows:

\begin{enumerate}[leftmargin=*]
    \item Construct the anchor prefix from the last $K = 90$ frames of the previously generated chunk, along with the corresponding audio context.
    \item Sample a latent variable $\mathbf{z}_n$ from the flow-matching prior (Section~\ref{sec:prior}), conditioned on speaker speech, identity, gesture type, and anchor-prefix.
    \item Decode the chunk using the temporal VAE decoder conditioned on $\mathbf{z}_n$, audio, and anchor-prefix.
    \item Apply overlap blending between consecutive chunks to smooth residual boundary artifacts.
\end{enumerate}

% -----------------------------------------------------------------------
\subsection{LLM-Based Gesture Type Planner}
\label{sec:planner}
% -----------------------------------------------------------------------
The VAE decoder requires a gesture type label $g$ to determine whether a segment should exhibit rhythmic beat gestures or meaning-driven semantic gestures.
While BEAT2 provides such annotations during training, these labels are not available for arbitrary speech at inference time.
To address this limitation, we introduce a gesture type planner that predicts structured segment-level annotations directly from speech transcripts.

\noindent \textbf{Task formulation.}
Given a word-level transcript with timestamps, the planner predicts a sequence of non-overlapping segments.
Each segment is defined by a start time, end time, and a gesture type:

\begin{verbatim}
[{"start_sec": 0.0, "end_sec": 2.4,
  "gesture_type": "beat", "text": "..."},
 ...]
\end{verbatim}
This segmentation output directly drives the chunk construction process in the autoregressive generation pipeline (Section~\ref{sec:autoreg}).

\noindent \textbf{Model and training.}
We finetune Qwen3.5-9B~\cite{qwen3} using LoRA~\cite{lora2022} and supervised finetuning.
Training data, ground-truth segmentations and gesture type labels, is derived from the BEAT2 dataset~\cite{emage2024}.
Each training sample is formatted as a chat-style interaction consisting of a system prompt defining the output schema, a transcript input with timestamps, and a structured JSON segment plan as the target output.
This formulation enables the model to jointly learn temporal segmentation and gesture type classification.

% \noindent \textbf{Data construction.}
% Since manual annotations are not available, BEAT2 \texttt{.sem} files are automatically converted into segment-level supervision.
% This allows the planner to learn dataset-specific gesture conventions without requiring additional human labeling.

\noindent \textbf{Few-shot and finetuned variants.}
We evaluate both a zero-/few-shot prompted Qwen3.5-9B and a LoRA finetuned variant trained with BEAT2-derived supervision.
The two models exhibit different behaviors: the few-shot variant produces conservative segmentations with fewer semantic gestures, while the finetuned model generates denser and more expressive segment structures.
A detailed comparison is provided in Section~\ref{sec:planner_eval}.

\noindent \textbf{Inference.}
At inference time, the planner defines the temporal segmentation and gesture type assignments used to construct chunks for the autoregressive decoder (Section~\ref{sec:autoreg}).

\subsection{Overall Training Strategy}
\label{sec:train}
GestureCraft is trained in a staged manner to decouple representation learning, temporal generation, and structured conditioning.
This design ensures stability while allowing each component to specialize before joint interaction.

\noindent \textbf{Training stages.}
We first train the hierarchical VAE described in Section~\ref{sec:vae} jointly with the anchor-conditioned objectives of Section~\ref{sec:autoreg}, using the reconstruction, KL, continuation velocity, and anchor-reconstruction losses.
The trained VAE encoder and decoder are then used to learn a conditional flow-matching latent prior (Section~\ref{sec:prior}) while keeping the VAE parameters fixed.
Finally, we optionally apply end-to-end perceptual finetuning of the prior to align generated motion with high-level quality objectives.

\noindent \textbf{High-level conditioning modules.}
The LLM-based gesture type planner (Section~\ref{sec:planner}) is trained independently and kept fixed during all motion model training stages.
Its predicted segmentations are used as conditioning signals for both the autoregressive decoder and the flow-matching prior, but are not updated.

\noindent \textbf{Frozen and trainable components.}
During flow-matching training, the VAE is frozen to provide a stable latent representation.
In the perceptual finetuning stage, the VAE decoder remains frozen and gradients are propagated to refine the flow-matching prior.

\noindent \textbf{Stage objectives.}
The system is not trained with a single joint loss.
Each stage optimizes its own objective:
(i) The VAE is trained with reconstruction and KL losses ($\mathcal{L}_\text{VAE}$),
(ii) The flow-matching prior is then trained with the rectified-flow loss ($\mathcal{L}_\text{flow}$) on top of the frozen VAE, and optionally finetuned with perceptual losses ($\mathcal{L}_\text{e2e}$), and
(iii) The anchor-conditioned decoder is trained jointly with the VAE under $\mathcal{L}_\text{cont}$ and the anchor-reconstruction loss.

% ======================================================================
% SECTION: Experiments
% ======================================================================

\section{Experiments}
\label{sec:experiments}
\noindent \textbf{Dataset.}
We use the BEAT2 (v2.0.0) English corpus~\cite{emage2024}, which contains recordings from 25 speakers with per-frame SMPL-X pose annotations, audio, transcripts, and gesture type labels.
We follow the official train/val/test split.
All prior methods that report results on BEAT2 evaluate only on Speaker~2 only.
However, we train a single model on all 25 speakers jointly and evaluate on both the Speaker~2 test set (for direct comparison with baselines) and the full 25-speakers test set (to verify cross-speaker generalization).

\noindent \textbf{Baselines.}
We compare against state-of-the-art co-speech gesture generation methods including recent VQ-based and diffusion-based approaches trained on BEAT2.
All baselines are evaluated under their reported settings when available, or re-evaluated using official checkpoints when provided.

% \noindent \textbf{Evaluation metrics.}
% We evaluate motion quality using Fr\'echet Gesture Distance (FGD) and Beat Consistency (BC).
% FGD measures distributional similarity between generated and ground-truth motion in a learned feature space, while BC measures alignment between motion dynamics and speech rhythm.

\noindent \textbf{Evaluation protocol.}
All hyperparameter and checkpoint selection decisions---including the classifier-free guidance scale $w$---are made on the validation split.
Final metrics are computed on the test split with 5 random seeds; we report mean~$\pm$~std.
% FGD and BC are reported in the scale used by prior work (raw metric $\times 10$).

\noindent \textbf{Implementation details.}
We optimize all components using Adam~\cite{adam2015} with learning rate $10^{-4}$ for VAE and decoder training, and $10^{-5}$ for flow-matching finetuning.
Models are trained for 400 epochs with batch size 8.
The KL coefficient is fixed at $\beta = 10^{-5}$ to prioritize reconstruction fidelity.

\subsection{Quantitative Comparison}
We report three standard metrics following \citet{emage2024}: Fr\'echet Gesture Distance~\cite{trimodal} (FGD, lower is better), Beat Consistency (BC, closer to ground-truth is better), and L1 diversity (L1Div, higher is better).
Our numbers are averaged over 5 random seeds; mean $\pm$ standard deviation are reported.
Baseline numbers in Table~\ref{tab:results} are sourced from \citet{syntalker}'s comparison table, in which all baselines are trained or finetuned only on Speaker~2.

\noindent \textbf{Speaker~2 comparison.}
Table~\ref{tab:results} compares our method against prior work on the
standard Speaker~2 test set.
Despite being trained on all 25 speakers simultaneously rather than finetuned on Speaker~2 alone, our model achieves the best FGD and the highest L1Div, with a BC within $0.05$ of ground-truth.

\noindent \textbf{Cross-speaker generalization.}
Table~\ref{tab:results_allspk} reports results on the full 25-speakers
test set.
We compare against EMAGE~\cite{emage2024} and SynTalker~\cite{syntalker}, both re-trained from their official codebases on the full 25-speakers training set with default hyperparameters.
Our model achieves FGD~$3.42$, roughly $6{\times}$ lower than SynTalker ($20.61$) and $7{\times}$ lower than EMAGE ($23.69$).
The EMAGE all-speakers model also collapses on Beat Consistency (BC~$0.85$ vs.\ ground-truth~$4.55$), indicating a broader failure to track speech rhythm at scale rather than an isolated FGD regression.
For reference, the released EMAGE checkpoint (trained only on Speaker~2) scores FGD~$25.1$ on this test set; re-training on all speakers yields only a modest improvement (FGD~$23.694$), indicating the limitation is architectural rather than a matter of training data coverage.
% To further verify that the aggregate result is not masking collapse on a subset of speakers, we evaluate the all-speakers prior (without any speaker-specific finetuning) on each speaker individually.
% Per-speaker FGD has median $8.86$ and ranges from $4.41$ on the easiest speaker to $34.1$ on the hardest, with $23$ of $25$ speakers below FGD $18$ and only one outlier above $20$.
% The hardest speakers correspond to those with the most distinctive motion styles in BEAT2.
% Per-speaker fine-tuning further reduces error on individual speakers, as demonstrated by the Speaker~2 results in Table~\ref{tab:results} ($\text{FGD} = 4.374$ for the spk2-fine-tuned model vs.\ $5.82$ for the all-speaker model evaluated on Speaker~2 alone).

\begin{table}[t]
\scriptsize
\caption{Quantitative comparison on BEAT2 Speaker~2 test set.
  % $\downarrow$ lower is better; $\uparrow$ higher is better;
  % $\rightarrow$ closer to GT is better.
  % All baselines are trained/fine-tuned on Speaker~2 only.
  % \textbf{Our model is trained on all 25 speakers} and evaluated here
  % on Speaker~2 for a fair comparison, demonstrating that
  % cross-speaker generalization does not sacrifice per-speaker quality.
  % Baseline numbers are sourced from \citet{syntalker}'s comparison
  % table; ``SynTalker (w/o both)'' refers to the SynTalker variant
  % trained without both motion-representation pre-training and
  % text--motion alignment pre-training, as defined in their paper.
  % The ground-truth row reports the canonical BEAT2 reference values
  % from \citet{emage2024}.
  % All metrics follow the standard BEAT2 reporting convention
  % (FGD$\times 10^{-1}$, BC$\times 10^{-1}$, diversity).
  Best in \textbf{bold} and second best \underline{underlined}.}
\label{tab:results}
\begin{tabular}{lccc}
\toprule
\textbf{Method}
  & \textbf{FGD $\downarrow$}
  & \textbf{BC $\rightarrow$}
  & \textbf{L1Div $\uparrow$} \\
\midrule
S2G~\cite{s2g}
  & 25.129           & \underline{6.902}  & 7.783  \\
Trimodal~\cite{trimodal}
  & 19.759           & 6.442  & 8.894  \\
HA2G~\cite{ha2g}
  & 19.364           & 6.601  & 9.671  \\
DisCo~\cite{disco}
  & 21.170           & 6.571  & 10.378 \\
Habibie~\cite{habibie}
  & 14.581           & 6.779  & 8.874  \\
CaMN~\cite{beat}
  & 8.752            & 6.731  & 9.279  \\
DiffuseStyleGesture~\cite{diffusestylegesture2023}
  & 10.137           & \textbf{6.891}  & 11.075 \\
TalkSHOW~\cite{talkshow2023}
  & 7.313            & 6.783  & 12.859 \\
EMAGE~\cite{emage2024}
  & 5.423            & 6.794  & 13.057 \\
SynTalker (w/o both)~\cite{syntalker}
  & \underline{4.687}& 7.363 & \underline{12.425} \\
\midrule
Ground Truth~\cite{emage2024}
  & 0.000            & 6.896             & 13.074            \\
\midrule
\textbf{GestureCraft (ours)}
  & \textbf{4.374}$^{\scriptstyle \pm 0.067}$
  & 6.944$^{\scriptstyle \pm 0.060}$
  & \textbf{14.85}$^{\scriptstyle \pm 0.250}$ \\
\bottomrule
\end{tabular}
\end{table}

% Convention: FGD/BC paper-scale = raw x 10 (evaluate_generation.py:418).
% Hands MPJPE rows in tab:recon_* are averaged across left/right hands.

\begin{table}[t]
\scriptsize
\caption{Quantitative comparison on BEAT2 full test set.
  % (265 recordings).
  % Both baselines are re-trained from the official codebase on the
  % full 25-speaker training set with default hyperparameters and the
  % $\text{\texttt{speaker\_dims}}\!=\!31$ multi-speaker setting,
  % selecting checkpoints by best validation FGD.
  % Baseline rows are single-seed; our row is mean over 5 seeds.
  % FGD and BC reported on the paper scale (raw~$\times 10$).
  }
\label{tab:results_allspk}
\begin{tabular}{lccc}
\toprule
\textbf{Method} & \textbf{FGD $\downarrow$} & \textbf{BC $\rightarrow$} & \textbf{L1Div $\uparrow$} \\
\midrule
EMAGE~\cite{emage2024}                 & 23.694            & 0.848               & 7.80              \\
SynTalker~\cite{syntalker}             & 20.605            & 6.154               & 8.51              \\
\textbf{GestureCraft (all-spk)}                & \textbf{3.416} & \textbf{4.348} & \textbf{9.311} \\
\midrule
Ground Truth                           & 0.000             & 4.555               & 8.427             \\
\bottomrule
\end{tabular}
\end{table}

\subsection{Ablation Studies}
We isolate the contribution of each component on the Speaker~2 test set (Table~\ref{tab:ablation}), using the same evaluation protocol as Table~\ref{tab:results}.

\begin{table}[t]
\scriptsize
\caption{Ablation on BEAT2 Speaker~2 test set.% (mean over 5 seeds, $w\!=\!0.5$).
  Each row removes or replaces one component of the full pipeline.}
\label{tab:ablation}
\begin{tabular}{lccc}
\toprule
\textbf{Configuration} & \textbf{FGD $\downarrow$} & \textbf{BC $\rightarrow$} & \textbf{L1Div $\uparrow$} \\
\midrule
Full model            & 4.374 & 6.944 & 14.85 \\
w/o anchor-prefix              & 5.183          & 6.721 & 14.77 \\
w/o flow-matching prior        & 7.384          & 7.681 & 9.94 \\
w/o E2E perceptual FT          & 5.373          & 6.791 & 15.10 \\
\midrule
Ground Truth                   & 0.000          & 6.896 & 13.074 \\
\bottomrule
\end{tabular}
\end{table}

\noindent \textbf{Anchor-prefix.}
Disabling anchor-prefix conditioning at inference degrades FGD by $0.81$ and pushes BC further from ground-truth.
This confirms that exposing the decoder to the previous chunk's motion context is necessary for temporal coherence across autoregressive boundaries.

\noindent \textbf{Flow-matching prior.}
Replacing the learned prior with $\mathbf{z} \!\sim\! \mathcal{N}(\mathbf{0}, \mathbf{I})$ sampling causes the largest single drop in quality, increasing FGD by $3.01$ and collapsing diversity (L1Div drops from $14.85$ to $9.94$).
BC counter-intuitively rises above ground truth, but this reflects over-animated, audio-uncorrelated motion that registers spurious beat hits rather than genuinely speech-aligned gestures.

\noindent \textbf{End-to-end perceptual finetuning.}
Skipping Stage~2 of the prior (using the all-speakers pretrained prior without speaker-specific perceptual finetuning) increases FGD by roughly $1.0$ on Speaker~2.
This isolates the contribution of the combined flow-matching, FGD moment-matching, and beat-alignment losses (Section~\ref{sec:prior}) from that of all-speakers pretraining alone.

\noindent \textbf{Continuation velocity loss.}
We isolate the continuation velocity loss by retraining the VAE with $\lambda_\text{cont} \!=\! 0$ (and the auxiliary anchor reconstruction weight set to zero), keeping all other settings identical.
Boundary metrics in Table~\ref{tab:cont_ablation} are computed in chunked autoregressive reconstruction mode (encode ground-truth chunks, decode each chunk using the previous decoded chunk as anchor), so the prior is not involved and the comparison isolates the decoder's boundary behavior.
Removing the loss roughly doubles the boundary velocity error (33.4 vs.~16.6\,mm) and the boundary jerk (37.9 vs.~17.3\,mm), confirming that the loss is the principal mechanism behind seamless chunk-to-chunk transitions; within-chunk reconstruction quality also improves slightly when the loss is enabled, suggesting it acts as a mild global regularizer in addition to its boundary-specific role.

\begin{table}[t]
\scriptsize
\caption{Continuation velocity loss ablation.% Speaker~2 test set,
  % chunked autoregressive reconstruction (encoder $\boldsymbol{\mu}$, no
  % prior). Boundary metrics computed across chunk transitions; within-chunk
  % MPJPE averaged across non-boundary frames. Lower is better.
  }
\label{tab:cont_ablation}
\begin{tabular}{lcc}
\toprule
\textbf{Metric (mm)} & \textbf{w/o $\mathcal{L}_\text{cont}$} & \textbf{w/ $\mathcal{L}_\text{cont}$} \\
\midrule
Boundary velocity error $\downarrow$ & 33.37 & \textbf{16.64} \\
Boundary jerk $\downarrow$           & 37.88 & \textbf{17.31} \\
Within-chunk MPJPE $\downarrow$      & 34.71 & \textbf{29.63} \\
\bottomrule
\end{tabular}
\end{table}

\noindent \textbf{LLM Planner Evaluation.}
\label{sec:planner_eval}
We evaluate two LLM planner variants on the Speaker~2 test set: a few-shot prompted Qwen3.5-9B and a LoRA fine-tuned
Qwen3.5-9B.
Both produce a per-segment beat vs. semantic segmentation, which replaces the ground-truth gesture type labels in the BEAT2 dataset during inference.
We report three categories of result:
(i)~Segmentation quality against ground-truth (Table~\ref{tab:planner_quality}),
(ii)~Downstream generation quality (Table~\ref{tab:planner_downstream}), and
(iii)~Impact of gesture type conditioning itself, isolated via all-beat and all-semantic ablations.
On Speaker~2, gesture type conditioning provides limited signal: the all-beat baseline ($4.337$) is within $0.04$ FGD of GT-conditioning ($4.374$), since $\sim$63\% of Speaker~2 chunks are beat.
The all-semantic ablation ($+1.55$ FGD) confirms that the conditioning channel is functional, but the all-beat result bounds how much any planner can gain from it on this distribution.
Both LLM planners cost $\sim$$0.2$ FGD relative to ground-truth labels (few-shot $+0.20$, LoRA $+0.26$).
The planner's contribution is therefore not a quality improvement but an engineering one: it removes the inference-time dependency on annotations at a small, well-characterized cost, demonstrating that transcripts alone supply enough signal to approximate BEAT2's gesture-type taxonomy.

\begin{table}[t]
\scriptsize
\caption{Segmentation quality of the LLM planners.}% on the Speaker~2
  % test set (15 recordings), measured against the BEAT2 \texttt{.sem}
  % ground truth.
  % Boundary~F1 uses a $0.5$\,s tolerance.}
\label{tab:planner_quality}
\begin{tabular}{lcc}
\toprule
\textbf{Metric} & \textbf{Few-shot} & \textbf{LoRA-FT} \\
\midrule
Per-frame accuracy $\uparrow$ & \textbf{75.0\%} & 57.8\% \\
Macro-F1 (beat + semantic) $\uparrow$ & \textbf{52.1\%} & 46.5\% \\
\quad Beat F1                 & \textbf{84.4\%} & 69.0\% \\
\quad Semantic F1             & 19.8\%          & \textbf{23.9\%} \\
Boundary F1 (@$0.5$\,s) $\uparrow$ & 19.1\% & \textbf{35.1\%} \\
\midrule
Predicted semantic ratio (GT $21.1\%$) & 11.3\% & 37.4\% \\
Mean segments per recording (GT $21.8$) & 12.5 & 48.8 \\
\bottomrule
\end{tabular}
\end{table}

\begin{table}[t]
\scriptsize
\caption{Generation quality based on gesture type source.}% (Speaker~2 test set, mean over 5 seeds, $w\!=\!0.5$).
  % All-beat / all-semantic isolate the contribution of gesture-type
  % conditioning itself.}
\label{tab:planner_downstream}
\begin{tabular}{lccc}
\toprule
\textbf{Label source} & \textbf{FGD $\downarrow$} & \textbf{BC $\rightarrow$} & \textbf{L1Div $\uparrow$} \\
\midrule
GT BEAT2 \texttt{.sem}      & \textbf{4.374} & 6.944 & 14.85 \\
LLM (few-shot)              & 4.575          & 6.906 & 14.69 \\
LLM (LoRA-FT)               & 4.634          & 7.132 & 15.15 \\
All-beat                    & 4.337          & 6.926 & 14.54 \\
All-semantic                & 5.882          & 7.019 & 15.99 \\
\midrule
Ground Truth                & 0.000          & 6.896 & 13.074 \\
\bottomrule
\end{tabular}
\end{table}


% \noindent \textbf{Implications.}

\noindent \textbf{VAE reconstruction upper bound vs.\ VQ baselines.}
Tables~\ref{tab:recon_spk2} and~\ref{tab:recon_allspk} report reconstruction quality using the encoder's posterior mean $\boldsymbol{\mu}$ (deterministic; no sampling noise) in chunked autoregressive mode.
We compare against the official pre-trained tokenizers of EMAGE and SynTalker (RVQ-VAE).
Our VAE and EMAGE's VQ-VAE share a chunked autoregressive evaluation protocol; SynTalker is evaluated under its own full-recording protocol (53 joints, zero global translation) since its non-chunked encoding makes a uniform protocol infeasible.
On the standard Speaker~2 test set, our continuous-latent VAE achieves the lowest joint error on both body and hands, with overall MPJPE about $1.7{\times}$ lower than the next-best VQ baseline (SynTalker) and $2.8{\times}$ lower than EMAGE; SynTalker's RVQ-VAE achieves lower FGD than our VAE on this set, but considerably higher joint error.
On the full 25-speakers test set the gap widens substantially: both VQ baselines were trained on Speaker~2 alone and degrade sharply outside their training distribution, while a single all-speakers instance of our VAE preserves reconstruction quality across the full corpus.
The all-speakers variant of our VAE remains within $\sim$2\,mm of the speaker-finetuned variant on Speaker~2, indicating that a single continuous tokenizer generalizes across the full corpus without per-speaker codebooks.

\begin{table}[t]
\scriptsize
\caption{Reconstruction upper bound on the BEAT2 Speaker~2 test set.}
  % (15 recordings; encoder $\boldsymbol{\mu}$, no sampling).
  % Per-part MPJPE in mm.
  % Our VAE and EMAGE's VQ-VAE are evaluated under a shared chunked
  % autoregressive protocol; SynTalker is evaluated using its own
  % protocol since their full-recording VQ encoding does not fit the
  % chunked autoregressive setup.
  % $^*$SynTalker: full-recording VQ encode/decode, 53 joints (excluding
  % eyes), zero global translation, per-recording averaging.}
\label{tab:recon_spk2}
\begin{tabular}{lcccc}
\toprule
 & \textbf{EMAGE} & \textbf{SynTalker$^*$} & \textbf{GestureCraft} & \textbf{GestureCraft} \\
 & \textbf{VQ-VAE} & \textbf{RVQ-VAE}      & \textbf{(all-spk)} & \textbf{(spk2-FT)}   \\
\midrule
FGD $\downarrow$         & 4.394 & \textbf{1.66} & 3.232 & 2.69 \\
MPJPE (all) $\downarrow$ & 90.55 & 55.72 & 34.84 & \textbf{32.66} \\
MPJPE body $\downarrow$  & 54.74 & 24.49 & 22.76 & \textbf{21.29} \\
MPJPE hands $\downarrow$ & 120.19 & 79.62 & 45.35 & \textbf{42.52} \\
\bottomrule
\end{tabular}
\end{table}

\begin{table}[t]
\scriptsize
\caption{Reconstruction upper bound on the full BEAT2 25-speakers test
  set.}% (265 recordings). Both VQ baselines are trained on Speaker~2 only
  % and degrade sharply outside that training distribution; our single
  % all-speaker continuous VAE preserves reconstruction quality across the
  % full corpus, with overall MPJPE roughly $4\times$ lower than either
  % VQ baseline.
  % $^*$SynTalker: full-recording VQ encode/decode, 53 joints (excluding
  % eyes), zero global translation.}
\label{tab:recon_allspk}
\begin{tabular}{lccc}
\toprule
 & \textbf{EMAGE} & \textbf{SynTalker$^*$} & \textbf{GestureCraft} \\
 & \textbf{VQ-VAE} & \textbf{RVQ-VAE} & \textbf{(all-spk)} \\
\midrule
FGD $\downarrow$         & 21.32  & 18.58  & \textbf{1.94} \\
MPJPE (all) $\downarrow$ & 112.25 & 113.66 & \textbf{26.03} \\
MPJPE body $\downarrow$  & 71.55  & 55.57  & \textbf{17.71} \\
MPJPE hands $\downarrow$ & 145.22 & 157.72 & \textbf{33.41} \\
\bottomrule
\end{tabular}
\end{table}

\noindent \textbf{Smoothness in continuous vs.\ discrete tokenizers.}
A central motivation for our continuous-latent VAE is that discrete codebooks introduce velocity discontinuities at codeword transitions.
Table~\ref{tab:smoothness} measures this directly: for each tokenizer, we run chunked autoregressive reconstruction of the Speaker~2 test set, forward through SMPL-X to obtain joint positions, and compute mean acceleration, mean jerk, and the 99th-percentile jerk (which captures boundary spikes invisible in mean statistics).
Our continuous VAE is the smoothest non-ground-truth method on every axis, with mean jerk $2.6{\times}$ lower than EMAGE.
The P99 metric is the most discriminative---it isolates the codeword transition artifacts---and quantifies the snapping pattern predicted by the discretization step in VQ tokenizers: our P99 jerk is $2.6{\times}$ GT while EMAGE's is $13.5{\times}$ GT.
SynTalker's RVQ sits between EMAGE and ours, consistent with multi-level residual quantization partially mitigating single-codebook
snapping.

\begin{table}[t]
\scriptsize
\caption{Joint-position smoothness on the Speaker~2 test set.}
  % (chunked autoregressive reconstruction, SMPL-X forward kinematics).
  % Lower is better; closer to ground truth is closer to natural motion.
  % P99 Jerk captures outlier spikes characteristic of codebook
  % transitions.}
\label{tab:smoothness}
\begin{tabular}{lccc}
\toprule
\textbf{Method} & \textbf{Mean Accel $\downarrow$} & \textbf{Mean Jerk $\downarrow$} & \textbf{P99 Jerk $\downarrow$} \\
\midrule
Ground Truth          & 3.70  & 2.33  & 16.73  \\
\textbf{Ours}         & \textbf{4.25}  & \textbf{5.33}  & \textbf{42.70}  \\
SynTalker RVQ-VAE     & 6.79  & 9.64  & 138.48 \\
EMAGE VQ-VAE          & 10.22 & 13.76 & 226.51 \\
\bottomrule
\end{tabular}
\end{table}

% ======================================================================
% SECTION: Conclusion and Limitations
% ======================================================================
\section{Conclusion}
\label{sec:conclusion}
We presented \textbf{GestureCraft}, a framework for audio-driven co-speech gesture synthesis built on a continuous, body-part-hierarchical latent rather than the discrete codebooks that dominate recent work.
Three design choices distinguish the system: a Gaussian latent partitioned channel-wise into face, body, and hand slices, each supervised by per-level body-part-localized partial-output losses; an anchor-prefix conditioning scheme paired with a continuation velocity loss that trains the decoder directly for kinematic continuity across chunk boundaries; and a conditional flow-matching prior pretrained on all $25$ BEAT2 speakers and fine-tuned end-to-end with FGD moment-matching and beat-alignment perceptual losses.
On the standard Speaker~2 benchmark, GestureCraft sets a new state of the art on FGD and diversity while keeping BC within $0.05$ of ground truth.
Trained jointly across all $25$ BEAT2 speakers, the system also provides, to our knowledge, the first published evaluation of cross-speaker generalization on the full BEAT2 25-speaker test set, outperforming the strongest re-trained VQ baselines on FGD by a wide margin.
% Ablations confirm that each component---anchor-prefix conditioning, the continuation velocity loss, the flow-matching prior, and the speaker-specific perceptual fine-tuning---contributes a measurable share of the final quality.
% A complementary LLM gesture-type planner removes the inference-time
% dependency on dataset-specific \texttt{.sem} annotations at a small, well-characterized cost ($\sim$$0.2$ FGD vs.\ GT labels), demonstrating that transcripts alone supply enough signal to approximate BEAT2's gesture-type taxonomy.

% \noindent \textbf{Limitations and future work.}
% Our evaluation is restricted to the BEAT2 corpus; we have not tested generalization to other gesture datasets or to in-the-wild speech, and gesture conventions may differ across languages and cultures beyond what BEAT2 covers.
% Generation runs offline through a $50$-step midpoint-method ODE integration of the flow-matching prior plus an autoregressive decoder; we have not characterized real-time performance.
% The body-part hierarchy is hand-designed (face / body / hands) rather than learned; whether a learned grouping would further improve reconstruction or generation remains an open question.
% The LLM planner's downstream impact on FGD is bounded by how much information gesture-type conditioning provides in the first place---we measure this directly through the \emph{all-beat} and \emph{all-semantic} ablations, but a stronger conditioning signal (for example, finer-grained or example-based) may unlock further gains.
% Finally, although our results include both quantitative comparison and per-component ablation, a comprehensive user study would further substantiate the perceptual claims; we leave it to future work.

% ======================================================================
% BIBLIOGRAPHY
% ======================================================================

% DO NOT INCLUDE ACKNOWLEDGMENTS IN AN ANONYMOUS SUBMISSION TO SIGGRAPH
%\begin{acks}
%\end{acks}

\bibliographystyle{ACM-Reference-Format}
\bibliography{sample-bibliography}

% ======================================================================
% FIGURE-ONLY PAGES (up to 2 pages allowed beyond the main text limit)
% ======================================================================
\clearpage

\begin{figure*}[p]
  \centering
  \includegraphics[width=\textwidth,height=0.78\textheight,keepaspectratio]{per_speaker_fgd.pdf}
  \Description{Three-panel figure showing per-speaker FGD on the full BEAT2
    25-speaker test set. (a) sorted bars; (b) pairwise scatter where all 25
    points lie above the y=x parity line; (c) per-speaker improvement ratio.}
  \caption{\textbf{Per-speaker FGD on the full BEAT2 25-speaker test set.}
    \textbf{(a)} Sorted horizontal bars (paper scale, lower is better);
    dashed lines mark each method's per-speaker mean.
    \textbf{(b)} Pairwise scatter: each point is one speaker; all $25$
    points lie above the diagonal $y\!=\!x$, so GestureCraft achieves lower
    FGD than EMAGE on every speaker.
    \textbf{(c)} Per-speaker improvement ratio (EMAGE\,/\,Ours):
    median $4.68{\times}$, range $1.68$--$8.20{\times}$.}
  \label{fig:per_speaker}
\end{figure*}

\clearpage

\begin{figure*}[p]
  \centering
  \includegraphics[width=\textwidth,height=0.82\textheight,keepaspectratio]{slice_ablation.pdf}
  \Description{Eight-row figure showing latent slice ablations on two
    BEAT2 Speaker~2 test clips. Top rows show ground truth; subsequent rows
    show the full-latent reconstruction and partial-output decoding with
    higher slices zeroed.}
  \caption{\textbf{Latent slice ablations on two BEAT2 Speaker~2 clips.}
    Top row: GT. Second row: full-latent reconstruction (encoder
    $\boldsymbol{\mu}$, no sampling), which closely tracks GT. Lower
    rows: partial-output decoding with higher slices zeroed.
    Reading from the bottom up: $\mathbf{z}^{(0)}$ alone already
    recovers most of the coarse pose; adding $\mathbf{z}^{(1)}$ makes the
    body motion richer and closer to GT; including $\mathbf{z}^{(2)}$
    (full $\mathbf{z}$) further improves fine hand and finger detail.
    The per-level loss supervises level~$\ell$ on body part $\ell$
    (face/body/hands), but the global loss on the full-latent output lets
    the decoder share information across slices, so ablation produces
    gracefully degraded motion rather than strict per-part isolation.}
  \label{fig:slice_ablation}
\end{figure*}